% =====================================================================
% CODEX ColombiaMacro · Parte II · Fundamentos para leer datos macro (español)
% Fragmento para \input: sin preámbulo.
% =====================================================================
\newcolumntype{L}[1]{>{\raggedright\arraybackslash}p{#1}}

\part{Fundamentos para leer datos macro}

Esta parte reúne las herramientas con las que se construye cada cifra del sitio. Cada concepto tiene una
definición, su fórmula, un ejemplo numérico y un recuadro de lectura o de cautela.

\begin{alerta}[Sobre los ejemplos]
Todos los números de los ejemplos de esta parte son \textbf{ilustrativos}: se eligieron para que las cuentas
sean fáciles de seguir y \emph{no} corresponden a datos actuales ni a ninguna fecha concreta. Las cifras
vigentes están siempre en el sitio.
\end{alerta}

Notación general: $X_t$ es el valor de una serie en el período $t$; $t-1$ es el período anterior (mes,
trimestre o día según la frecuencia), $t-4$ es el mismo trimestre del año anterior y $t-12$ el mismo mes del año
anterior. Las tasas de crecimiento se expresan en porcentaje: $g=100\,(X_t/X_{t-k}-1)$.

% ---------------------------------------------------------------------
\chapter{Niveles, cambios y unidades}
\label{es:f:cambios}

\section{Niveles y cambios}
Un \textbf{nivel} es el valor de una variable en un momento (la TRM de hoy, el PIB de un trimestre, la tasa de
desempleo de un mes). Un \textbf{cambio} compara dos niveles. Casi todo lo que interesa en macroeconomía es un
cambio: cuánto creció el PIB, cuánto subieron los precios, cuánto se movió la tasa de un bono.

Hay dos familias de variables, y cada una cambia en una unidad distinta:
\begin{itemize}
\item \textbf{Niveles}, medidos en unidades (pesos, dólares, puntos de un índice, personas). Su cambio se mide
  en \textbf{porcentaje}: $\Delta\%X = 100\,(X_t/X_{t-k}-1)$.
\item \textbf{Tasas}, medidas ya en porcentaje (inflación, desempleo, tasas de interés, crecimiento, razones al
  PIB). Su cambio se mide en \textbf{puntos porcentuales}: $\Delta X = X_t-X_{t-k}$.
\end{itemize}

\section{Porcentaje, puntos porcentuales y puntos básicos}
\label{es:f:pp}
\begin{itemize}
\item Un \textbf{punto porcentual} (pp) es la diferencia aritmética entre dos porcentajes.
\item Un \textbf{punto básico} (pb) es una centésima de punto porcentual: $1\pp = 100$ pb. Se usa para tasas de
  interés y de cambio que se mueven en fracciones pequeñas.
\item El \textbf{cambio relativo} de una tasa es otra cosa: $100\,(X_t/X_{t-k}-1)$.
\end{itemize}

\begin{dato}[Ejemplo ilustrativo]
La tasa de desempleo pasa de 10,0\% a 11,0\%. El cambio es $+1{,}0$ pp ($+100$ pb). En términos relativos,
el desempleo subió $100\,(11/10-1)=10\%$. Una tasa de interés que pasa de 4\% a 5\% sube 1 pp, pero un 25\% en
términos relativos.
\end{dato}

\begin{alerta}
Confundir pp con \% es el error más común al leer datos macro. El sitio siempre expresa en pp los cambios de las
tasas y en \% los de los niveles (TRM, COLCAP, tasa de cambio real, valores en dólares). Cuando lea ``la inflación
subió 1,1'', verifique si son puntos (de 5,1\% a 6,2\%) o por ciento (de 5,1\% a 5,16\%).
\end{alerta}

\section{Variación mensual}
\[
m_t = 100\left(\frac{X_t}{X_{t-1}}-1\right).
\]
Es la medida más reciente posible, pero también la más ruidosa: en datos sin desestacionalizar refleja sobre todo
la temporada (matrículas en febrero, alimentos con el clima, primas en diciembre).

\begin{dato}[Ejemplo ilustrativo]
Un índice de precios pasa de 120,0 a 120,6: $m_t=100\,(120{,}6/120{,}0-1)=0{,}50\%$.
\end{dato}

\section{Variación anual}
\[
a_t = 100\left(\frac{X_t}{X_{t-12}}-1\right)\ \text{(mensual)},\qquad
a_t = 100\left(\frac{X_t}{X_{t-4}}-1\right)\ \text{(trimestral)}.
\]
Compara con el mismo período del año anterior, por lo que elimina la mayor parte de la estacionalidad sin necesidad
de modelos. Es la medida principal del sitio para la inflación, el PIB y el empleo. Con variaciones mensuales,
\[
1+\frac{a_t}{100}=\prod_{k=0}^{11}\left(1+\frac{m_{t-k}}{100}\right).
\]

\begin{dato}[Ejemplo ilustrativo]
Si el índice vale 126,0 hoy y 120,0 hace doce meses, la variación anual es $100\,(126/120-1)=5{,}0\%$. Si los
precios subieran 0,4\% todos los meses, la inflación anual sería $100\,(1{,}004^{12}-1)=4{,}91\%$, no $12\times0{,}4=4{,}8\%$.
\end{dato}

\begin{leer}
La variación anual es lenta: resume doce meses. Un cambio de tendencia aparece primero en las variaciones
mensuales o trimestrales anualizadas y solo después en la anual.
\end{leer}

\section{Variación acumulada en el año (año corrido)}
Compara con el cierre del año anterior (diciembre o el cuarto trimestre):
\[
c_t = 100\left(\frac{X_t}{X_{\text{dic},\,a-1}}-1\right).
\]
Para flujos (PIB, exportaciones) el \emph{año corrido} compara la suma de los períodos transcurridos con la
misma suma del año anterior:
\[
c_t = 100\left(\frac{\sum_{q=1}^{Q}Y_{a,q}}{\sum_{q=1}^{Q}Y_{a-1,q}}-1\right).
\]

\begin{dato}[Ejemplo ilustrativo]
Si el IPC de diciembre fue 122,0 y el de junio es 125,05, la inflación en lo corrido del año es $2{,}5\%$. Si el PIB
del primer semestre suma 520 (miles de millones de pesos constantes) frente a 505 un año antes, el crecimiento
año corrido es $100\,(520/505-1)=2{,}97\%$.
\end{dato}

\begin{alerta}
La variación año corrido de enero a junio no es comparable con una variación anual: cubre menos meses y hereda la
estacionalidad del primer semestre.
\end{alerta}

\section{Variación trimestral anualizada}
\label{es:f:anualizada}
Para ver el ritmo reciente se usa el cambio frente al trimestre anterior de la serie
\emph{desestacionalizada} ($Y^{sa}$), llevado a un año por capitalización compuesta:
\[
g^{\text{anualizada}}_t = 100\left[\left(\frac{Y^{sa}_t}{Y^{sa}_{t-1}}\right)^{4}-1\right]
=100\left[(1+g^{q}_t)^4-1\right],
\]
donde $g^q_t$ es la variación trimestral en tanto por uno. Es la convención de la Oficina de Análisis Económico
de Estados Unidos (BEA). En datos mensuales, el sitio usa la versión de tres meses contra tres meses:
\[
g^{3m/3m}_t = 100\left[\left(\frac{\bar x_{t,3}}{\bar x_{t-3,3}}\right)^{4}-1\right],\qquad
\bar x_{t,3}=\tfrac13(x_t+x_{t-1}+x_{t-2}).
\]

\begin{dato}[Ejemplo ilustrativo]
Si el PIB desestacionalizado crece 1,3\% en un trimestre, el ritmo anualizado es $100\,(1{,}013^4-1)=5{,}3\%$:
``si la economía creciera así durante un año entero, crecería 5,3\%''. Si el promedio de los últimos tres meses del
ISE es 101,2 y el de los tres anteriores 100,0, el ritmo 3m/3m anualizado es $100\,(1{,}012^4-1)=4{,}9\%$.
\end{dato}

\begin{leer}
Cuando el ritmo anualizado está por encima de la variación anual, la actividad \emph{acelera}; cuando está por
debajo, \emph{se modera}. La página \emph{Ciclo} usa esta comparación.
\end{leer}

\begin{alerta}
La anualización multiplica el ruido por cuatro: un trimestre atípico (un paro, un festivo, una revisión) produce
saltos grandes. Léala junto a la variación anual, nunca sola.
\end{alerta}

\section{Sumas y promedios móviles de 12 meses}
\label{es:f:moviles}
Para flujos mensuales con mucha estacionalidad (exportaciones, importaciones, remesas, balanza cambiaria, matrículas
de empresas) el sitio usa la \textbf{suma móvil de 12 meses}:
\[
S^{12}_t=\sum_{k=0}^{11}x_{t-k},\qquad \Delta\% S^{12}_t = 100\left(\frac{S^{12}_t}{S^{12}_{t-12}}-1\right).
\]
Para tasas (desempleo por sexo, subutilización) usa el \textbf{promedio móvil}, $\bar x^{12}_t=S^{12}_t/12$; y para
series trimestrales, la suma de los últimos \textbf{cuatro trimestres}. La GEIH publica además \textbf{trimestres
móviles}: el dato de un mes es el promedio de ese mes y los dos anteriores (por ejemplo, may--jul).

\begin{dato}[Ejemplo ilustrativo]
Si las exportaciones de los últimos 12 meses suman US\$52.000 millones y las de los 12 meses previos US\$48.000
millones, el crecimiento en 12 meses es $100\,(52/48-1)=8{,}3\%$.
\end{dato}

\begin{alerta}
Dos trimestres móviles consecutivos comparten dos de sus tres meses: no son observaciones independientes. Una
suma de 12 meses reacciona con lentitud: un cambio brusco tarda un año en reflejarse por completo.
\end{alerta}

\section{Efecto base}
\label{es:f:base}
La variación anual de hoy depende tanto del mes que entra como del mes que \emph{sale} de la ventana de doce
meses:
\[
1+\pi^{12}_t=\left(1+\pi^{12}_{t-1}\right)\frac{1+m_t}{1+m_{t-12}}.
\]
Si hace un año hubo un salto atípico ($m_{t-12}$ alto), la inflación anual cae cuando ese mes sale de la ventana,
aunque los precios de hoy se comporten igual. Lo mismo ocurre al revés.

\begin{dato}[Ejemplo ilustrativo]
La inflación anual del mes pasado era 6,0\%. Este mes los precios suben 0,3\%, pero en el mismo mes del año
anterior habían subido 2,0\% por un choque puntual. La nueva inflación anual es
$100\,[1{,}06\times1{,}003/1{,}02-1]=4{,}23\%$: cae 1,77 pp por efecto base, no porque los precios hayan
bajado.
\end{dato}

\begin{leer}
Antes de interpretar un cambio grande en una variación anual, mire la variación mensual de hoy \emph{y} la del
mismo mes del año anterior. Si la segunda explica el movimiento, es efecto base. El crecimiento del PIB en 2021,
medido contra el trimestre más bajo de 2020, es el ejemplo extremo.
\end{leer}

\section{Estacionalidad y series desestacionalizadas}
\label{es:f:estacional}
Muchas series tienen un patrón que se repite cada año (\textbf{estacionalidad}) y efectos de calendario (número
de días hábiles, Semana Santa, festivos). Hay tres formas de neutralizarlos:
\begin{enumerate}
\item \textbf{Comparar con el mismo período del año anterior} (variación anual): la más simple y la usada en las
  series \emph{originales}.
\item \textbf{Usar la serie desestacionalizada} ($sa$) que publica la fuente (por ejemplo el PIB, el ISE y el
  desempleo del DANE), estimada con modelos estadísticos que remueven la estacionalidad y el calendario. Solo
  con ella tiene sentido comparar un período con el inmediatamente anterior.
\item \textbf{Sumas o promedios de 12 meses}, que contienen exactamente una temporada de cada tipo.
\end{enumerate}

\begin{alerta}
No compare un mes o trimestre con el anterior en una serie original: el cambio estará dominado por la temporada.
Y recuerde que las series desestacionalizadas se \emph{revisan}: cada dato nuevo cambia ligeramente la
estimación de los factores estacionales de la historia reciente.
\end{alerta}

% ---------------------------------------------------------------------
\chapter{Precios, valores reales e índices}
\label{es:f:precios}

\section{Nominal y real}
Una magnitud \textbf{nominal} se mide a precios corrientes (pesos de cada año); una \textbf{real} se mide a precios
constantes (pesos de un año base) o como volumen. La relación es
\[
\text{Real}_t=\frac{\text{Nominal}_t}{P_t/100},
\qquad 1+g^{\text{real}}=\frac{1+g^{\text{nominal}}}{1+\pi},
\]
con $P_t$ un índice de precios (base 100) y $\pi$ su variación. El PIB del DANE se publica en las dos versiones;
el crecimiento de la economía es siempre el \emph{real}.

\begin{dato}[Ejemplo ilustrativo]
Si el PIB nominal crece 9,0\% y los precios de lo producido suben 5,5\%, el crecimiento real es
$100\,(1{,}09/1{,}055-1)=3{,}3\%$.
\end{dato}

\begin{alerta}[Ilusión nominal]
Un salario, una utilidad o un saldo de crédito que crece 8\% con una inflación de 9\% \emph{cae} en términos
reales. El sitio deflacta por la inflación las variables donde esa diferencia importa (crédito, tasas de
interés, PIB).
\end{alerta}

\section{Deflactores}
\label{es:f:deflactor}
Un \textbf{deflactor} es el índice de precios con el que se pasa de nominal a real.
\begin{itemize}
\item \textbf{IPC}: precios de la canasta de consumo de los hogares (DANE). Es el que define la inflación y la
  meta del Banco de la República.
\item \textbf{Deflactor implícito del PIB}: precios de todo lo que se \emph{produce} en el país,
  \[
  D_t = 100\,\frac{\text{PIB nominal}_t}{\text{PIB real}_t}.
  \]
  Incluye exportaciones (petróleo, carbón, café) y excluye importaciones; por eso se separa del IPC cuando
  cambian los términos de intercambio (Kohli, 2004).
\end{itemize}

\begin{leer}
Si el deflactor del PIB sube más que el IPC, los precios de lo que el país vende (incluidas sus exportaciones)
suben más que los de lo que consume: suele coincidir con mejores términos de intercambio. Si sube menos, ocurre
lo contrario.
\end{leer}

\section{Inflación: total, de fondo y aportes}
\label{es:f:inflacion}
La \textbf{inflación} es la variación anual del IPC: $\pi_t=100\,(IPC_t/IPC_{t-12}-1)$. La meta del Banco de la
República es 3\% con un rango de $\pm1$ pp.

La \textbf{inflación de fondo} (o básica) excluye los componentes más volátiles para medir la presión persistente.
La principal del sitio es la \emph{inflación sin alimentos ni regulados} (cálculo de BanRep). Si la total sube y
la de fondo no, el choque probablemente es transitorio (clima, energía, tarifas).

El \textbf{aporte} de un componente $j$ a la inflación total (en pp) combina su peso en la canasta y su
variación; de forma aproximada,
\[
\text{aporte}_{j,t}\approx w_{j,t-12}\,\pi_{j,t},
\]
donde $w_{j,t-12}$ es el peso efectivo del componente al inicio de la ventana (ponderación de la canasta ajustada
por los precios relativos). El DANE publica los aportes exactos, que suman la inflación total.

\begin{dato}[Ejemplo ilustrativo]
Si alimentos pesa 0,15 de la canasta y su inflación anual es 10\%, aporta cerca de $0{,}15\times10=1{,}5$ pp a la
inflación total. Un arriendo que pesa 0,25 con inflación de 5\% aporta $1{,}25$ pp: un componente grande con
inflación moderada puede aportar tanto como uno pequeño con inflación alta.
\end{dato}

\section{Tasa de interés real: Fisher exacta y aproximada}
\label{es:f:fisher}
La \textbf{ecuación de Fisher} relaciona la tasa nominal $i$, la tasa real $r$ y la inflación $\pi$ (en tanto
por uno):
\[
1+i=(1+r)(1+\pi)\quad\Longrightarrow\quad r=\frac{1+i}{1+\pi}-1\qquad\text{(exacta)},
\]
\[
r\approx i-\pi\qquad\text{(aproximada)}.
\]
La diferencia entre ambas es $r_{\text{aprox}}-r=\dfrac{\pi\,(i-\pi)}{1+\pi}$: pequeña con tasas bajas y
relevante con tasas de dos dígitos. El sitio usa siempre la forma exacta.

\begin{itemize}
\item \textbf{Tasa real ex ante}: deflacta con la inflación \emph{esperada} (en el sitio, la implícita a 1 año
  de los TES). Es la que importa para decidir hoy un crédito o una inversión.
\item \textbf{Tasa real ex post}: deflacta con la inflación \emph{observada}. En el sitio se usa la última
  inflación publicada en cada fecha (el IPC de un mes se considera disponible el día 10 del mes siguiente), para
  no usar información que aún no existía.
\end{itemize}

\begin{dato}[Ejemplo ilustrativo]
Con una tasa nominal de 12\% y una inflación esperada de 7\%: $r=1{,}12/1{,}07-1=4{,}67\%$. La aproximación daría
5,0\%, un error de 0,33 pp.
\end{dato}

\begin{leer}
Una tasa nominal alta no implica una política restrictiva si la inflación esperada también es alta. Lo que frena o
estimula la economía es la tasa \emph{real} frente a la neutral (sección~\ref{es:f:neutral}).
\end{leer}

\section{Índices base 100 y rebase}
\label{es:f:indices}
Un \textbf{índice} expresa una serie en relación con un período base:
\[
I_t=100\,\frac{X_t}{X_{b}}.
\]
Un valor de 115 significa 15\% por encima del período base. Para cambiar la base a otro período $t_0$
(\textbf{rebase}) se divide toda la serie por su valor en $t_0$:
\[
\tilde I_t = 100\,\frac{I_t}{I_{t_0}}.
\]
El rebase no altera las variaciones porcentuales entre dos fechas; solo cambia el punto de referencia. Por eso
permite comparar series de unidades distintas (acciones, monedas, componentes de la inversión) desde un mismo
punto de partida.

\begin{dato}[Ejemplo ilustrativo]
Una serie vale 80, 92 y 110 en tres fechas. Con base en la primera: 100, 115 y 137,5. El valor final (137,5) dice
directamente que la serie subió 37,5\% en el período.
\end{dato}

\begin{alerta}
En un gráfico re-basado, el orden y la distancia entre líneas dependen del punto de partida. Dos series que parten de
100 en una fecha favorable para una de ellas mostrarán una ventaja que desaparece con otro punto de partida. Pruebe
varios horizontes.
\end{alerta}

\section{Logaritmos multiplicados por 100}
\label{es:f:log}
Para variaciones pequeñas, la diferencia de logaritmos por 100 es casi igual a la variación porcentual:
\[
100\,[\ln X_t-\ln X_{t-1}]\approx100\left(\frac{X_t}{X_{t-1}}-1\right).
\]
Tiene dos ventajas: las variaciones en logaritmos se \textbf{suman} en el tiempo y son simétricas (subir y bajar
la misma cantidad logarítmica devuelve al punto inicial). Por eso la brecha del producto y los filtros de tendencia
trabajan con $y_t=100\ln Y_t$; una brecha de $g$ significa que el PIB está aproximadamente $g\%$ por encima de su
tendencia. La conversión exacta es $100\,(e^{g/100}-1)$.

\begin{dato}[Ejemplo ilustrativo]
Si $X$ pasa de 100 a 105, la variación es 5,00\% y la diferencia logarítmica $100\ln1{,}05=4{,}88$. Para una
brecha de $+0{,}3$ la diferencia entre ambas medidas es despreciable.
\end{dato}

% ---------------------------------------------------------------------
\chapter{Composición: participaciones, aportes y concentración}
\label{es:f:composicion}

\section{Participaciones}
La \textbf{participación} de un componente $i$ en un total es $w_{i,t}=100\,X_{i,t}/X_t$. Las participaciones
suman 100 cuando los componentes son aditivos (precios corrientes, conteos de personas, valores en dólares).

\section{Aportes al crecimiento}
\label{es:f:aportes}
El \textbf{aporte} (o contribución) de un componente al crecimiento anual del total, en pp, es
\[
c_{i,t}=100\,\frac{X_{i,t}-X_{i,t-4}}{X_{t-4}}=\frac{w_{i,t-4}\;g_{i,t}}{100},
\]
es decir, el crecimiento del componente ponderado por su participación \emph{un año antes}. Si los componentes son
aditivos, $\sum_i c_{i,t}=g_t$: los aportes suman el crecimiento total.

\begin{dato}[Ejemplo ilustrativo]
Un sector que pesaba 30\% del valor agregado y crece 2\% aporta $0{,}30\times2=0{,}6$ pp. Otro que pesaba 2\% y
crece 15\% aporta $0{,}02\times15=0{,}3$ pp. El sector grande que crece poco aporta el doble que el pequeño que crece
mucho.
\end{dato}

\begin{leer}
Para saber \emph{quién explica} el crecimiento, mire los aportes, no las tasas de crecimiento. Para saber
\emph{quién está acelerando}, compare la tasa de crecimiento de cada componente con su propia historia.
\end{leer}

\section{Por qué los volúmenes encadenados no suman exacto}
\label{es:f:encadenados}
El DANE mide el PIB real con \textbf{índices de volumen encadenados} (referencia 2015): cada año se valora con los
precios del año anterior y los eslabones se encadenan. Este método refleja mejor los cambios de precios relativos,
pero tiene una consecuencia: los componentes en volúmenes encadenados \textbf{no son aditivos}. La suma de los 12
sectores no es exactamente el valor agregado total, las participaciones calculadas con volúmenes pueden sumar algo
distinto de 100 y los aportes suman \emph{aproximadamente} el crecimiento.

Además, el PIB por el lado de la oferta es el valor agregado de los sectores \emph{más los impuestos netos de
subvenciones}; por eso los aportes sectoriales suman el crecimiento del valor agregado, no el del PIB. Por el lado
de la demanda, el sitio muestra la diferencia como \emph{Existencias y discrepancia}, que incluye la variación de
inventarios y la no aditividad.

\begin{alerta}
Una diferencia pequeña entre la suma de aportes y el crecimiento total no es un error de cálculo: es una propiedad
del encadenamiento y del tratamiento de los impuestos.
\end{alerta}

\section{Crecimiento sin un sector}
Para saber cuánto crece la economía sin un sector $X$ (por ejemplo, sin \emph{Gobierno, educación y salud}):
\[
g^{\,\text{sin }X}_t=100\,\frac{\sum_i c_{i,t}-c_{X,t}}{\sum_i w_{i,t-4}-w_{X,t-4}}.
\]

\begin{dato}[Ejemplo ilustrativo]
Los aportes suman 3,0 pp con pesos que suman 100. El sector $X$ pesa 18 y aporta 1,8 pp. Sin él, el resto crece
$100\times(3{,}0-1{,}8)/(100-18)=1{,}46\%$: más de la mitad del crecimiento venía de un sector que pesa menos de una
quinta parte.
\end{dato}

\section{Índices de difusión y amplitud}
\label{es:f:difusion}
Un \textbf{índice de difusión} cuenta cuántos componentes se mueven en una dirección (Burns y Mitchell, 1946). El
sitio usa dos:
\begin{itemize}
\item \emph{Sectores que crecen, de 12}: número de sectores con variación anual positiva.
\item \emph{Amplitud del ciclo}: número de sectores que producen por encima de su propia tendencia. Más de 6 de 12
  se lee como expansión \emph{generalizada}; 6 o menos, como expansión \emph{concentrada}.
\end{itemize}

\begin{leer}
Un crecimiento total alto con pocos sectores que crecen es frágil: depende de pocos motores. Una inflación alta con
la mayoría de subclases subiendo por encima de la meta (distribución desplazada a la derecha) es más persistente que
una explicada por pocos rubros.
\end{leer}

\section{Herfindahl y número equivalente}
\label{es:f:herfindahl}
El \textbf{índice de Herfindahl--Hirschman} mide la concentración con la suma de las participaciones al cuadrado
(en tanto por uno):
\[
H=\sum_{i=1}^{n}s_i^2,\qquad \frac1n\le H\le 1,\qquad N^{\text{eq}}=\frac1H .
\]
El \textbf{número equivalente} $N^{\text{eq}}$ es el número de componentes de igual tamaño que darían la misma
concentración. El sitio lo usa para la diversificación de la estructura productiva de cada departamento, de los
destinos de exportación y de los grupos de productos exportados.

\begin{dato}[Ejemplo ilustrativo]
Con participaciones de 40\%, 30\%, 20\% y 10\%: $H=0{,}16+0{,}09+0{,}04+0{,}01=0{,}30$ y $N^{\text{eq}}=3{,}3$. Si
los cuatro pesaran 25\%, $H=0{,}25$ y $N^{\text{eq}}=4$.
\end{dato}

\begin{alerta}
El número equivalente depende del nivel de detalle de la clasificación: con cinco grupos de productos, su máximo
posible es 5. Solo compare números equivalentes calculados con la misma clasificación.
\end{alerta}

% ---------------------------------------------------------------------
\chapter{Tendencia, ciclo y capacidad}
\label{es:f:ciclo}

\section{Tendencia y ciclo}
Una serie de actividad se puede descomponer en una \textbf{tendencia} (la trayectoria de largo plazo, asociada a
la capacidad productiva) y un \textbf{ciclo} (las desviaciones transitorias alrededor de ella):
\[
y_t=\tau_t+c_t,\qquad y_t=100\ln Y_t .
\]
Ni la tendencia ni el ciclo se observan: se \emph{estiman}. Distintos métodos dan resultados distintos, sobre todo
en los extremos de la muestra. Por eso el sitio muestra varios y su rango.

\section{El filtro de Hodrick y Prescott (HP)}
\label{es:f:hp}
La tendencia HP (Hodrick y Prescott, 1997) es la serie $\tau$ que resuelve
\[
\min_{\{\tau_t\}}\ \sum_{t=1}^{T}(y_t-\tau_t)^2+\lambda\sum_{t=2}^{T-1}\big[(\tau_{t+1}-\tau_t)-(\tau_t-\tau_{t-1})\big]^2 .
\]
El primer término pide que la tendencia se parezca a los datos; el segundo penaliza los cambios en su pendiente. El
parámetro $\lambda$ regula la suavidad: $\lambda=1.600$ para datos trimestrales y $\lambda=129.600$ para mensuales
(equivalencia de Ravn y Uhlig, 2002). En forma matricial, $\tau=(I+\lambda D^\top D)^{-1}y$, con $D$ la matriz de
segundas diferencias.

\subsection{Tiempo real frente a dos colas}
\begin{itemize}
\item \textbf{HP de dos colas} ($\tau_{t|T}$): usa toda la muestra, incluidas observaciones posteriores a $t$. Es
  la mejor estimación \emph{ex post}, pero la tendencia de un año cambia cuando llegan datos nuevos.
\item \textbf{HP en tiempo real} o de una cola ($\tau_{t|t}$): para cada $t$ se resuelve el problema solo con
  $y_1,\dots,y_t$ y se conserva el último punto. Reproduce lo que se habría estimado en cada fecha. Es la
  estimación principal del sitio (con un mínimo de 20 trimestres).
\end{itemize}
Orphanides y van Norden (2002) muestran que las revisiones entre ambas versiones son del mismo orden que la brecha
misma: la estimación del último trimestre es la más incierta.

\subsection{El tratamiento de la pandemia}
Los trimestres 2020-T2 a 2021-T2 se reemplazan por una interpolación lineal \emph{solo para estimar la tendencia}
(en datos mensuales, marzo de 2020 a junio de 2021). La brecha se mide siempre contra el dato observado. Sin este
tratamiento, una caída de casi una quinta parte del PIB en un trimestre arrastraría la tendencia y deformaría la
lectura de los años vecinos.

\section{Otros métodos: Hamilton, Christiano--Fitzgerald y Beveridge--Nelson}
\label{es:f:metodos}
\begin{itemize}
\item \textbf{Hamilton (2018)}: el ciclo es el error de proyectar el PIB ocho trimestres adelante con sus cuatro
  últimos valores, estimado por mínimos cuadrados:
  \[
  y_{t+8}=\beta_0+\sum_{k=0}^{3}\beta_k\,y_{t-k}+v_{t+8},\qquad c^{H}_{t+8}=\hat v_{t+8}.
  \]
  Evita las propiedades espurias del HP en series con raíz unitaria.
\item \textbf{Christiano y Fitzgerald (2003)}: filtro de paso de banda que conserva las fluctuaciones con períodos de
  6 a 32 trimestres (1,5 a 8 años), la duración típica de un ciclo.
\item \textbf{Beveridge y Nelson (1981)}: la tendencia es el nivel al que convergería el PIB si no hubiera nuevos
  choques, con un AR(4) para el crecimiento.
\end{itemize}
La página \emph{Ciclo} muestra las cinco brechas (HP en tiempo real, HP de dos colas, Hamilton, CF y BN), su
\textbf{mediana} como consenso, su \textbf{rango} como medida de incertidumbre y cuántas son positivas.

\begin{leer}
Si los cinco métodos coinciden en el signo, la lectura (por encima o por debajo de la capacidad) es robusta. Si se
reparten, la economía está cerca de su capacidad y la fase es incierta: conviene esperar el siguiente dato.
\end{leer}

\section{Brecha del producto y capacidad}
\label{es:f:brecha}
La \textbf{brecha del producto} es la distancia porcentual entre el PIB observado y su tendencia (la
\textbf{capacidad} o PIB potencial):
\[
g_t=y_t-\tau_t=100\,(\ln Y_t-\ln Y^*_t)\quad\Longleftrightarrow\quad Y^*_t=Y_t\,e^{-g_t/100}.
\]
Una brecha positiva indica que la economía produce por encima de lo que puede sostener sin presiones de inflación
(empresas a tope, mercado laboral estrecho); una negativa indica \textbf{holgura} (capacidad ociosa, desempleo). La
capacidad es como una \emph{frontera de posibilidades de producción} que se desplaza con el tiempo.

\begin{dato}[Ejemplo ilustrativo]
Si el PIB trimestral desestacionalizado es 260 (billones de pesos de 2015) y la brecha es $+0{,}3\%$, la capacidad
es $260\,e^{-0{,}003}=259{,}2$: la economía produce unos 0,8 billones por encima de su capacidad en el trimestre.
\end{dato}

\begin{alerta}
La brecha es una \emph{estimación}: su nivel exacto cambia con el método y con los datos nuevos. Lo más informativo
es su signo cuando los métodos coinciden y su \emph{dirección} (si se abre o se cierra).
\end{alerta}

\section{Crecimiento potencial (``ritmo habitual'')}
\label{es:f:potencial}
El \textbf{crecimiento potencial} es el crecimiento anual de la tendencia:
\[
g^{\text{pot}}_t=100\left[\exp\!\left(\frac{\tau_t-\tau_{t-4}}{100}\right)-1\right].
\]
El sitio lo calcula con la tendencia HP de dos colas (la versión en tiempo real oscila con el rebote posterior a la
pandemia) y lo presenta como el \emph{ritmo habitual} de la economía. Crecer por encima de él cierra una brecha
negativa o abre una positiva; crecer por debajo hace lo contrario.

\begin{leer}
Crecimiento y capacidad son cosas distintas. Una economía que crece 3,5\% puede estar \emph{por debajo} de su
capacidad si venía de muy abajo, y una que crece 1\% puede estar por encima. Lo relevante para la inflación es la
distancia a la capacidad, no la tasa de crecimiento.
\end{leer}

\section{El reloj del ciclo}
\label{es:f:reloj}
La fase del ciclo combina el \textbf{nivel} de la brecha y su \textbf{dirección}, medida como su cambio en dos
trimestres, $\Delta_2 g_t=g_t-g_{t-2}$ (dos trimestres reducen el ruido de un solo dato):

\begin{center}
\begin{tabular}{@{}lcc@{}}
\toprule
 & $\Delta_2 g_t\ge 0$ & $\Delta_2 g_t<0$\\\midrule
$g_t\ge 0$ & \textcolor{cmgreen}{\textbf{Expansión}} & \textcolor{cmamber}{\textbf{Desaceleración}}\\
$g_t<0$ & \textcolor{cmblue}{\textbf{Recuperación}} & \textcolor{cmred}{\textbf{Contracción}}\\
\bottomrule
\end{tabular}
\end{center}

Es la lógica del \emph{business cycle clock} de la OCDE. En datos mensuales (ISE) la dirección es el cambio en tres
meses de la brecha mensual.

\begin{dato}[Ejemplo ilustrativo]
Si la brecha es $-0{,}4\%$ y hace dos trimestres era $-1{,}0\%$, entonces $\Delta_2 g=+0{,}6$ pp: la economía está
por debajo de su capacidad pero acercándose, en \emph{Recuperación}. Si en el trimestre siguiente la brecha pasa a
$+0{,}2\%$ y sigue subiendo, entra en \emph{Expansión}.
\end{dato}

\begin{leer}
La economía suele girar en el sentido del reloj: expansión $\rightarrow$ desaceleración $\rightarrow$ contracción
$\rightarrow$ recuperación. Cerca de cero, un movimiento pequeño cambia el cuadrante: lea el recorrido de varios
trimestres, no un punto.
\end{leer}

\section{Ley de Okun}
\label{es:f:okun}
La ley de Okun (1962) relaciona la brecha del producto con la brecha del desempleo:
\[
u_t-u^*_t=\beta\,g_t+\varepsilon_t,
\]
donde $u^*_t$ es la tendencia del desempleo (HP) y $\beta<0$. El sitio estima $\beta$ por mínimos cuadrados sin
2020--2021. Un $\beta$ cercano a cero indica que el empleo responde poco a la producción (en Colombia, en parte
porque la informalidad absorbe los choques).

\begin{dato}[Ejemplo ilustrativo]
Con $\beta=-0{,}2$, una brecha del producto de $+1\%$ se asocia, en promedio, con un desempleo 0,2 pp por debajo de
su tendencia.
\end{dato}

\section{Fechado de expansiones y recesiones}
\label{es:f:bryboschan}
El sitio identifica \textbf{picos} y \textbf{valles} en el nivel del ISE (promedio móvil de tres meses) con la regla de
Bry y Boschan (1971): un pico es un máximo local en una ventana de $\pm6$ meses, un valle un mínimo local, y picos y
valles deben alternarse (Harding y Pagan, 2002). Una \textbf{recesión} va de un pico a un valle; una
\textbf{expansión}, de un valle al pico siguiente. Para cada fase se reporta su duración en meses y la variación
del ISE.

\begin{alerta}
Es un fechado académico y descriptivo, no oficial. El último giro siempre es provisional: un pico solo se confirma
cuando han pasado suficientes meses de caída.
\end{alerta}

% ---------------------------------------------------------------------
\chapter{Herramientas estadísticas}
\label{es:f:estadistica}

\section{Percentil histórico}
\label{es:f:percentil}
El \textbf{percentil histórico} ubica el último dato $x^*$ dentro de su propia historia desde 2010:
\[
p=100\times\frac{\#\{\,t\ge 2010:\ x_t\le x^*\,\}}{\#\{\,t\ge2010\,\}}.
\]
Un percentil de 96 significa que el dato actual es mayor o igual que el 96\% de las observaciones desde 2010. El
sitio lo traduce a palabras: muy bajo ($<10$), bajo ($<30$), normal (30--70), alto ($\le90$) y muy alto ($>90$).

\begin{dato}[Ejemplo ilustrativo]
Si una serie tiene 200 observaciones desde 2010 y el último dato es mayor o igual que 160 de ellas, su percentil es
80: ``Alto''.
\end{dato}

\begin{alerta}
El percentil depende de la ventana: desde 2010 incluye la pandemia y el ciclo de inflación posterior. No dice si un
nivel es bueno o malo, ni si es sostenible; solo qué tan inusual es frente a esa historia.
\end{alerta}

\section{Volatilidad anualizada}
\label{es:f:volatilidad}
La \textbf{volatilidad} es la desviación estándar de los cambios diarios en una ventana móvil, llevada a un año con
la raíz del número de días hábiles (252):
\[
\sigma^{\text{anual}}_t=\sqrt{252}\times\operatorname{desv.est.}\{\Delta x_{t-59},\dots,\Delta x_t\}.
\]
Para la tasa de cambio, $\Delta x$ son cambios logarítmicos diarios (en \%); para los TES, cambios diarios de la
tasa (en pb). El sitio usa ventanas de 60 días hábiles. La regla $\sqrt{252}$ supone cambios diarios
independientes.

\begin{dato}[Ejemplo ilustrativo]
Si la tasa del TES a 10 años cambia, en promedio, con una desviación estándar de 6 pb diarios, su volatilidad
anualizada es $6\sqrt{252}\approx95$ pb. Si la TRM se mueve con una desviación de 0,7\% diaria, su volatilidad anual
es cerca de 11\%.
\end{dato}

\begin{leer}
La volatilidad mide el tamaño de los movimientos, no su dirección. Compárela con su promedio histórico (línea
punteada) y con la de las monedas vecinas: una volatilidad alta solo en Colombia apunta a factores locales.
\end{leer}

\section{Correlación}
\label{es:f:correlacion}
El coeficiente de correlación de Pearson entre dos series de cambios $a$ y $b$ en una ventana de $n$ observaciones es
\[
\rho_{ab}=\frac{\sum_t(a_t-\bar a)(b_t-\bar b)}{\sqrt{\sum_t(a_t-\bar a)^2}\sqrt{\sum_t(b_t-\bar b)^2}},\qquad -1\le\rho\le1 .
\]
El sitio la calcula en ventanas móviles de 52 semanas con variaciones semanales porcentuales (por ejemplo, TRM
frente al petróleo Brent y frente al índice amplio del dólar).

\begin{alerta}[Correlación no es causalidad]
Una correlación alta dice que dos variables se movieron juntas en la ventana, no que una cause a la otra, ni que la
relación vaya a mantenerse. Las correlaciones móviles cambian de signo y de tamaño con el tiempo.
\end{alerta}

\section{Rectas de ajuste descriptivas}
Algunos gráficos de dispersión incluyen una recta de mínimos cuadrados, $y=a+bx$. Describe la asociación promedio en
el período, no un modelo causal ni una proyección. El coeficiente $b$ dice cuánto cambia $y$, en promedio, por cada
unidad de $x$.

% ---------------------------------------------------------------------
\chapter{Mercado laboral}
\label{es:f:laboral}

\section{Tasas básicas}
Con PET la población en edad de trabajar (15 años y más), FT la fuerza de trabajo (ocupados más desocupados),
$O$ los ocupados y $D$ los desocupados:
\[
TGP=100\,\frac{FT}{PET},\qquad TO=100\,\frac{O}{PET},\qquad TD=100\,\frac{D}{FT},\qquad
TD=100\left(1-\frac{TO}{TGP}\right).
\]

\begin{dato}[Ejemplo ilustrativo]
Con PET de 40 millones, FT de 26 millones y 23,4 millones de ocupados: $TGP=65\%$, $TO=58{,}5\%$ y
$TD=100\,(1-58{,}5/65)=10\%$.
\end{dato}

\begin{leer}
El desempleo puede bajar por dos razones muy distintas: porque se crean empleos (sube la TO) o porque la gente deja
de buscar (baja la TGP). Mire siempre las tres tasas juntas.
\end{leer}

\section{Subutilización}
Siguiendo la 19.ª CIET de la OIT (2013), además del desempleo el sitio muestra la \textbf{tasa combinada de
desocupación y subocupación por horas} y la \textbf{medida compuesta de subutilización}, que agrega la fuerza de
trabajo potencial (personas que quieren trabajar pero no buscan o no están disponibles):
\[
TCSD=100\,\frac{D+SH}{FT},\qquad MCSFT=100\,\frac{D+SH+FTP}{FT+FTP}.
\]

\section{Informalidad}
La \textbf{tasa de informalidad} es la proporción de ocupados que trabajan sin seguridad social o en unidades
económicas no registradas (definición del DANE de 2023, alineada con la OIT). Por sector o ciudad,
$\text{tasa}=100\times\text{informales}/\text{ocupados}$.

\begin{alerta}
La serie vigente empieza en 2021 y no es comparable con la serie anterior (otra definición y marco muestral). La
informalidad cambia despacio: medio punto en un año es un movimiento relevante.
\end{alerta}

% ---------------------------------------------------------------------
\chapter{Tasas de interés y curva de rendimientos}
\label{es:f:tasas}

\section{Tasas efectivas anuales}
Las tasas del sitio son \textbf{efectivas anuales} (e.a.): ya incluyen la capitalización dentro del año. Una tasa
nominal $j$ capitalizable $m$ veces al año equivale a una efectiva $i=(1+j/m)^m-1$. Comparar tasas con distinta
convención introduce errores de varias décimas cuando las tasas son altas.

\section{Tasa de política, tasa neutral y postura}
\label{es:f:neutral}
La \textbf{tasa de política monetaria} (TPM) es la tasa mínima a la que el Banco de la República presta liquidez a
un día a las entidades financieras. La \textbf{tasa real neutral} es la tasa real que ni estimula ni frena la
economía (Laubach y Williams, 2003); no se observa y se estima con incertidumbre. La \textbf{postura} compara la tasa
real ex ante con la neutral:
\[
\text{brecha de postura}=r^{\text{ex ante}}-r^{n}.
\]
Por encima de la neutral la política es restrictiva; por debajo, expansiva. El rango y el margen que usa el sitio
están en la sección~\ref{es:i:estados-tabla}.

\section{Traspaso (transmisión) de la tasa de política}
\label{es:f:traspaso}
El \textbf{traspaso} mide qué parte del cambio de la TPM en un ciclo llegó a otra tasa $k$ (IBR, CDT, crédito):
\[
\text{traspaso}_k=100\times\frac{i_k(t_1)-i_k(t_0)}{TPM(t_1)-TPM(t_0)},
\]
con $t_0$ el día anterior a la primera decisión del ciclo y $t_1$ tres meses después de la última. Un traspaso de
100\% significa que la tasa se movió lo mismo que la TPM.

\begin{dato}[Ejemplo ilustrativo]
Si en un ciclo de alzas la TPM sube 4,0 pp y la tasa de consumo 2,6 pp, el traspaso es 65\%.
\end{dato}

\section{Crecimiento real del crédito}
\[
g^{\text{real}}=100\left[\frac{1+g^{\text{nominal}}/100}{1+\pi/100}-1\right].
\]
\begin{dato}[Ejemplo ilustrativo]
Si la cartera crece 8\% nominal con una inflación de 5\%, el crédito real crece $100\,(1{,}08/1{,}05-1)=2{,}9\%$.
\end{dato}

\section{La curva cero cupón}
\label{es:f:curva}
La \textbf{curva de rendimientos} muestra la tasa que exige el mercado por prestarle al Gobierno a cada plazo. Una
tasa \textbf{cero cupón} $y_h$ es la de un bono que paga todo al vencimiento (sin cupones intermedios); su precio
por cada peso a recibir en $h$ años es
\[
P_h=\frac{1}{(1+y_h)^h}.
\]
El Banco de la República estima cada día las curvas cero cupón de los TES con el modelo de Nelson y Siegel (1987),
en pesos (tasa nominal) y en UVR (tasa real, porque el principal se ajusta con la inflación), y publica los plazos de
1, 5 y 10 años.

\section{Nivel, pendiente, curvatura y forma}
\label{es:f:factores}
Con las tasas en pesos a 1, 5 y 10 años:
\[
\text{nivel}=\frac{y_1+y_5+y_{10}}{3},\qquad \text{pendiente}=y_{10}-y_1,\qquad
\text{curvatura}=2y_5-y_1-y_{10}.
\]
Estos tres factores resumen casi todo el movimiento de una curva (Litterman y Scheinkman, 1991). La pendiente define la
\textbf{forma}: \emph{normal} (creciente), \emph{plana} o \emph{invertida} (los plazos cortos pagan más que los
largos).

\begin{dato}[Ejemplo ilustrativo]
Con $y_1=9{,}0\%$, $y_5=9{,}6\%$ e $y_{10}=10{,}2\%$: nivel 9,6\%, pendiente $+1{,}2$ pp y curvatura
$2(9{,}6)-9{,}0-10{,}2=0$ (curva recta). Si $y_5$ fuera 9,9\%, la curvatura sería $+0{,}6$: la parte media se arquea
hacia arriba.
\end{dato}

\begin{leer}
Una pendiente positiva es lo normal: prestar a más plazo exige compensación por el tiempo y el riesgo. Una curva plana
o invertida suele aparecer cuando la tasa de política está alta y el mercado espera que baje; la pendiente ha tenido
contenido informativo sobre la actividad futura en muchos países (Estrella y Hardouvelis, 1991), pero el sitio no la
usa para pronosticar.
\end{leer}

\section{Tipos de movimiento de la curva}
Comparando dos fechas, con $\Delta N$ el cambio del nivel y $\Delta S$ el de la pendiente (umbral de 0,10 pp):
\begin{center}\small
\begin{tabular}{@{}lll@{}}
\toprule
 & $\Delta S\ge0$ (se empina) & $\Delta S<0$ (se aplana)\\\midrule
$\Delta N\ge0$ (suben las tasas) & Alza con empinamiento (\emph{bear steepening}) & Alza con aplanamiento (\emph{bear flattening})\\
$\Delta N<0$ (bajan las tasas) & Baja con empinamiento (\emph{bull steepening}) & Baja con aplanamiento (\emph{bull flattening})\\
\bottomrule
\end{tabular}
\end{center}
Si $|\Delta N|$ y $|\Delta S|$ son menores que 0,10 pp, el sitio dice \emph{sin cambio de fondo}. ``Bear'' y ``bull''
se refieren al precio de los bonos: cuando las tasas suben, los precios bajan.

\section{Inflación implícita (breakeven)}
\label{es:f:breakeven}
Para un plazo $h$, con $i_h$ la tasa TES en pesos y $r_h$ la tasa TES en UVR:
\[
\pi^{BE}_h=100\left[\frac{1+i_h}{1+r_h}-1\right].
\]
Es la inflación promedio que iguala el rendimiento de ambos bonos: si la inflación resulta mayor, gana el bono en UVR;
si resulta menor, gana el bono en pesos. Por la identidad de Fisher, el cambio de la tasa nominal se reparte entre el
cambio de la tasa real y el de la compensación por inflación:
\[
\Delta i_h\approx\Delta r_h+\Delta\pi^{BE}_h .
\]

\begin{dato}[Ejemplo ilustrativo]
Con $i_5=10{,}0\%$ y $r_5=4{,}5\%$: $\pi^{BE}_5=100\,(1{,}100/1{,}045-1)=5{,}26\%$.
\end{dato}

\begin{alerta}[No es una expectativa pura]
La inflación implícita incluye una \textbf{prima por riesgo inflacionario} (compensación por la incertidumbre sobre
la inflación) y una \textbf{prima por liquidez} (los TES UVR se negocian menos), y la UVR se indexa con rezago a la
inflación pasada, lo que afecta sobre todo al plazo de 1 año (Gürkaynak, Sack y Wright, 2010). Por eso el sitio la
describe como \emph{lo que el mercado descuenta} y como una cota de la expectativa pura.
\end{alerta}

\section{Tasas forward, la 5y5y y la trayectoria a 10 años}
\label{es:f:forward}
La \textbf{tasa forward} entre los años $a$ y $b$ es la tasa anual implícita para ese tramo futuro, que hace
equivalente invertir a $b$ años o a $a$ años y reinvertir:
\[
f_{a,b}=\left[\frac{(1+y_b)^b}{(1+y_a)^a}\right]^{\frac{1}{b-a}}-1 .
\]
La \textbf{5y5y} es la inflación implícita promedio entre los años 5 y 10:
\[
f^{nom}_{5,10}=\left[\frac{(1+i_{10})^{10}}{(1+i_5)^{5}}\right]^{1/5}\!-1,\quad
f^{real}_{5,10}=\left[\frac{(1+r_{10})^{10}}{(1+r_5)^{5}}\right]^{1/5}\!-1,\quad
\pi^{5y5y}=\frac{1+f^{nom}_{5,10}}{1+f^{real}_{5,10}}-1 .
\]
Es el indicador estándar de \textbf{anclaje} de las expectativas: a ese plazo, la política monetaria debería haber
llevado la inflación a la meta. Con las implícitas cero cupón $b_1$, $b_5$ y $b_{10}$, la \textbf{trayectoria} que
descuenta el mercado tiene tres tramos que encadenan exactamente la implícita a 10 años:
\[
\text{año 1: } b_1,\qquad f_{1,5}=\left[\frac{(1+b_5)^5}{1+b_1}\right]^{1/4}\!-1,\qquad
f_{5,10}=\left[\frac{(1+b_{10})^{10}}{(1+b_5)^5}\right]^{1/5}\!-1,
\]
\[
(1+b_1)(1+f_{1,5})^4(1+f_{5,10})^5=(1+b_{10})^{10}.
\]

\begin{dato}[Ejemplo ilustrativo]
Con $i_5=10{,}0\%$, $i_{10}=10{,}5\%$, $r_5=4{,}5\%$ y $r_{10}=4{,}8\%$: $f^{nom}=11{,}00\%$, $f^{real}=5{,}10\%$ y
$\pi^{5y5y}=1{,}1100/1{,}0510-1=5{,}61\%$. Con $b_1=5{,}5\%$, $b_5=5{,}0\%$ y $b_{10}=4{,}8\%$, la trayectoria es 5,5\% el
primer año, $f_{1,5}=4{,}88\%$ por año entre los años 1 y 5 y $f_{5,10}=4{,}60\%$ por año entre los años 5 y 10: el
mercado descuenta una inflación que baja, pero que no llega al 3\%.
\end{dato}

\begin{leer}
Si la 5y5y se mantiene dentro del rango meta, el mercado confía en que la inflación volverá a la meta; por encima del
techo de 4\% se considera que las expectativas están \emph{desancladas}. Cada punto de la historia de la 5y5y es lo que
el mercado esperaba \emph{ese día} para los años 5 a 10: no es un pronóstico de su propia línea.
\end{leer}

\section{Primas sobre la tasa de política}
La diferencia entre una tasa de mercado y la TPM vigente (por ejemplo $y_{10}-TPM$) combina lo que el mercado espera de la
TPM en el futuro y una \textbf{prima por plazo} (compensación por mantener un bono largo, que crece y se vuelve más
volátil con el vencimiento). El sitio la muestra como dato, sin descomponerla.

% ---------------------------------------------------------------------
\chapter{Tasa de cambio y activos}
\label{es:f:cambio}

\section{TRM, apreciación y depreciación}
\label{es:f:trm}
La \textbf{TRM} (tasa representativa del mercado) es el número de pesos por un dólar, certificado cada día por la
Superintendencia Financiera con las operaciones del día hábil anterior. Como se expresa en pesos por dólar:
\begin{itemize}
\item si la TRM \textbf{sube}, el peso se \textbf{deprecia} (se debilita): se necesitan más pesos por dólar;
\item si la TRM \textbf{baja}, el peso se \textbf{aprecia} (se fortalece).
\end{itemize}
El cambio del dólar en pesos y el del peso en dólares no son simétricos:
\[
\Delta\%(\text{COP por USD})=100\left(\frac{E_t}{E_{t-k}}-1\right),\qquad
\Delta\%(\text{USD por COP})=100\left(\frac{E_{t-k}}{E_t}-1\right).
\]

\begin{dato}[Ejemplo ilustrativo]
Si la TRM pasa de 4.000 a 4.400, el dólar sube 10\% en pesos y el peso cae $100\,(4.000/4.400-1)=-9{,}1\%$ en dólares.
\end{dato}

\section{Tasas cruzadas}
Para expresar el peso frente a otra moneda $M$ se combinan la TRM y las unidades de $M$ por dólar:
\[
\text{COP por }M=\frac{\text{COP por USD}}{M\text{ por USD}} .
\]
\begin{dato}[Ejemplo ilustrativo]
Con TRM de 4.000 y 5,0 reales por dólar, un real cuesta $4.000/5{,}0=800$ pesos.
\end{dato}

\begin{leer}[¿Es el peso o es el dólar?]
Si la TRM sube y el dólar global (índice amplio de la Reserva Federal) también sube con fuerza, el movimiento es en
buena parte del dólar. Si el peso se mueve más que el real brasileño, el peso mexicano o el sol peruano, hay factores
propios de Colombia.
\end{leer}

\section{Tasa de cambio real (ITCR)}
\label{es:f:itcr}
La \textbf{tasa de cambio real} ajusta la tasa nominal por los precios relativos. Para un socio $j$, con $E_j$ los
pesos por unidad de su moneda, $P_j$ sus precios y $P$ los de Colombia,
\[
Q_j=\frac{E_j\,P_j}{P}.
\]
El ITCR multilateral del Banco de la República promedia los socios comerciales con ponderaciones de comercio
(forma geométrica, base 2010 = 100):
\[
ITCR_t=100\prod_j\left(\frac{Q_{j,t}}{Q_{j,2010}}\right)^{w_j}.
\]
Un \textbf{aumento} del ITCR es una \textbf{depreciación real} (los bienes colombianos se abaratan frente a los
extranjeros); una disminución, una apreciación real.

\begin{dato}[Ejemplo ilustrativo]
Si en un año el peso se deprecia 5\% frente a un socio, la inflación del socio es 3\% y la de Colombia 6\%, la tasa real
bilateral cambia $100\,[1{,}05\times1{,}03/1{,}06-1]=2{,}0\%$: depreciación real de 2\%, menor que la nominal porque
la inflación colombiana fue mayor.
\end{dato}

\begin{alerta}
El año base (2010) es solo una fecha de comparación, no un nivel de equilibrio. Que el ITCR esté por debajo de su
promedio histórico no significa que ``deba'' subir.
\end{alerta}

\section{Rendimiento en dólares}
\label{es:f:usd}
Para un inversionista que mide en dólares, el rendimiento de un activo en pesos combina el precio local y la tasa de
cambio:
\[
1+R^{USD}=\frac{1+R^{COP}}{1+\Delta\%E/100},
\]
con $\Delta\%E$ la variación de la TRM. El COLCAP en dólares es $\text{COLCAP}_t/E_t$.

\begin{dato}[Ejemplo ilustrativo]
Si el COLCAP sube 20\% y la TRM baja 10\% (el peso se aprecia), el rendimiento en dólares es
$100\,(1{,}20/0{,}90-1)=33{,}3\%$.
\end{dato}

\section{Índices bursátiles}
\label{es:f:bolsa}
\begin{itemize}
\item \textbf{Ponderado por capitalización} (como el COLCAP): cada acción pesa según su capitalización ajustada; las
  empresas grandes mueven el índice.
\item \textbf{Equiponderado}: todas las acciones pesan igual, con rebalanceo periódico. Con $R_{j,t}$ el retorno
  semanal de la acción $j$ y $A_t$ las acciones disponibles,
  \[
  I_t=I_{t-1}\Big(1+\frac{1}{|A_t|}\sum_{j\in A_t}R_{j,t}\Big),\qquad I_0=100 .
  \]
\item \textbf{Índice de precios frente a retorno total}: el COLCAP y los índices construidos en el sitio no incluyen
  dividendos.
\item \textbf{Retorno a 12 meses}: último dato frente al último disponible 365 días antes.
\item \textbf{Caída desde el máximo} (\emph{drawdown}): $100\,(P_t/\max_{s\le t}P_s-1)$.
\end{itemize}

\begin{leer}
Si el índice ponderado sube más que el equiponderado, las ganancias se concentran en las empresas grandes; si ocurre
lo contrario, el alza es amplia. Una bolsa que sube por pocas empresas es más sensible a un choque en una de ellas
(Gabaix, 2011).
\end{leer}

\begin{alerta}[Sesgo de supervivencia]
Aplicar la canasta \emph{actual} hacia atrás omite las empresas que salieron del índice (por adquisiciones,
deslistamientos o malos resultados) y tiende a sobrestimar el rendimiento pasado.
\end{alerta}

% ---------------------------------------------------------------------
\chapter{Cuentas externas, fiscales y comercio}
\label{es:f:externo}

\section{FOB y CIF}
\label{es:f:fobcif}
\begin{itemize}
\item \textbf{FOB} (\emph{free on board}): valor de la mercancía puesta en el puerto de salida, sin fletes ni seguros
  internacionales. Las exportaciones se valoran FOB.
\item \textbf{CIF} (\emph{cost, insurance and freight}): valor de la mercancía más el flete y el seguro hasta el
  destino. Las importaciones del DANE se valoran CIF:
  \[
  \text{CIF}=\text{FOB}+\text{flete}+\text{seguro}.
  \]
\end{itemize}

\begin{dato}[Ejemplo ilustrativo]
Una importación de US\$100 FOB con US\$6 de flete y US\$1 de seguro vale US\$107 CIF. Una balanza comercial
calculada como exportaciones FOB menos importaciones CIF es por eso algo más negativa que la de la balanza de pagos,
que valora ambos lados FOB.
\end{dato}

\section{Valor, precio y volumen}
El cambio del valor de las exportaciones o importaciones se puede separar en precio y volumen implícito:
\[
1+g^{\text{valor}}=(1+g^{\text{precio}})(1+g^{\text{volumen}})\quad\Longrightarrow\quad
g^{\text{volumen}}=\frac{1+g^{\text{valor}}}{1+g^{\text{precio}}}-1 .
\]
Los \textbf{términos de intercambio} son el cociente entre el índice de precios de exportación y el de importación.

\begin{dato}[Ejemplo ilustrativo]
Si el valor exportado sube 8\% y los precios de exportación 5\%, el volumen implícito sube $100\,(1{,}08/1{,}05-1)=2{,}9\%$.
\end{dato}

\begin{alerta}
El volumen implícito es un \emph{residuo}, no una medición directa de cantidades: hereda los errores del índice de
precios.
\end{alerta}

\section{Razones al PIB y conversión a dólares}
\label{es:f:razones}
Para comparar flujos de tamaño muy distinto se dividen por el PIB. En el sitio:
\begin{itemize}
\item el numerador y el denominador cubren el mismo período (normalmente la suma de los \textbf{últimos cuatro
  trimestres});
\item cuando el numerador está en dólares (cuenta corriente, exportaciones, remesas, posición de inversión), el PIB
  nominal en pesos se convierte a dólares con la TRM promedio de cada trimestre:
  \[
  PIB^{US\$}=\sum_{q=1}^{4}\frac{N_q}{\overline{E}_q},\qquad
  \text{razón}=100\times\frac{\sum_{q=1}^{4}X^{US\$}_q}{PIB^{US\$}} .
  \]
\end{itemize}
La \textbf{apertura comercial} es $(X+M)/PIB$.

\begin{dato}[Ejemplo ilustrativo]
Si el PIB de cuatro trimestres equivale a US\$450.000 millones y el déficit de cuenta corriente de esos trimestres es
US\$13.500 millones, el déficit es $100\times13{,}5/450=3{,}0\%$ del PIB.
\end{dato}

\begin{alerta}
Una razón al PIB puede cambiar porque cambia el numerador, el PIB o la tasa de cambio. Una depreciación fuerte reduce
el PIB en dólares y \emph{sube} las razones de las variables en dólares aunque estas no cambien.
\end{alerta}

\section{Balanza de pagos}
\label{es:f:bp}
La balanza de pagos registra las transacciones de los residentes con el resto del mundo (Manual del FMI, sexta
edición, MBP6). Sus cuentas son:
\begin{itemize}
\item \textbf{Cuenta corriente} (CC): bienes, servicios, \emph{ingreso primario} (utilidades, dividendos e
  intereses) e \emph{ingreso secundario} (transferencias, sobre todo remesas);
\item \textbf{Cuenta de capital} (CK): transferencias de capital (pequeña en Colombia);
\item \textbf{Cuenta financiera} (CF): adquisición neta de activos financieros menos emisión neta de pasivos:
  inversión directa, de cartera, otra inversión (préstamos) y activos de reserva.
\end{itemize}
La identidad contable es
\[
CC+CK=CF+\varepsilon,
\]
donde $\varepsilon$ son los \textbf{errores y omisiones} netos, con el signo que cierra la identidad. Con los signos del
MBP6, una CF \emph{negativa} significa que el país se endeuda o recibe inversión neta del exterior (los pasivos crecen
más que los activos); una CF positiva, que presta o invierte afuera. Un déficit de cuenta corriente se financia con una
CF negativa.

\begin{leer}[Convención de los gráficos del sitio]
Para facilitar la lectura, el gráfico de financiamiento muestra la cuenta financiera \textbf{con el signo
invertido}: positivo = entrada neta de capital. En esa presentación, las reservas positivas indican una
\emph{disminución} de reservas (una fuente de financiamiento) y las negativas, una acumulación.
\end{leer}

\begin{dato}[Ejemplo ilustrativo]
Con una cuenta corriente de $-3{,}5\%$ del PIB y una cuenta de capital nula, la cuenta financiera es aproximadamente
$-3{,}5\%$. Presentada con signo invertido, podría verse como $+3{,}0$ de inversión directa, $+0{,}8$ de cartera,
$+0{,}2$ de otra inversión y $-0{,}5$ de reservas (acumulación): suman $+3{,}5$.
\end{dato}

\begin{leer}
Importa tanto el tamaño del déficit como su financiamiento: un déficit financiado con inversión extranjera directa es más
estable que uno financiado con flujos de cartera de corto plazo (Calvo, Leiderman y Reinhart, 1993).
\end{leer}

\section{Deuda externa y posición de inversión internacional}
La \textbf{deuda externa} son los pasivos con no residentes que exigen pagos de principal o intereses (préstamos,
bonos, créditos comerciales), pública y privada. La \textbf{posición de inversión internacional} (PII) es el saldo de
todos los activos y pasivos financieros externos; su posición neta es $\text{PII}^{neta}=\text{Activos}-\text{Pasivos}$
(Lane y Milesi-Ferretti, 2007). Los déficits de cuenta corriente acumulados hacen más negativa la PII; también la
mueven los cambios de precios y de tasa de cambio.

\section{Balance fiscal total y primario}
\label{es:f:fiscal}
Para el Gobierno nacional central (cifras de caja, suma de cuatro trimestres sobre el PIB de esos trimestres):
\[
\text{balance total}=\text{ingresos}-\text{gastos},\qquad
\text{balance primario}=\text{balance total}+\text{intereses}.
\]
El balance primario dice si el Gobierno cubre sus gastos sin contar el costo de la deuda. La razón
$\text{intereses}/\text{ingresos}$ mide la carga financiera.

\begin{dato}[Ejemplo ilustrativo]
Con ingresos de 16\% del PIB y gastos de 22\% (de los cuales 4,5\% son intereses): balance total $-6{,}0\%$ y balance
primario $-1{,}5\%$. Los intereses absorben $4{,}5/16=28\%$ de los ingresos.
\end{dato}

\section{Dinámica de la deuda pública}
Con $d_t$ la deuda sobre el PIB, $i$ la tasa de interés nominal implícita, $\gamma$ el crecimiento del PIB nominal y
$pb_t$ el balance primario (en \% del PIB):
\[
d_t=d_{t-1}\,\frac{1+i}{1+\gamma}-pb_t\quad(\text{sin otros flujos que generan deuda}).
\]
Si $i>\gamma$, se necesita un superávit primario para estabilizar la deuda; si $i<\gamma$, la deuda puede estabilizarse con
un déficit primario moderado (Blanchard, 2019). La deuda en moneda extranjera agrega el efecto de la tasa de cambio.

\begin{dato}[Ejemplo ilustrativo]
Con $d_{t-1}=60\%$, $i=8\%$, $\gamma=9\%$ y un déficit primario de 1,5\% ($pb=-1{,}5$):
$d_t=60\times1{,}08/1{,}09+1{,}5=61{,}0\%$. La deuda sube aunque el crecimiento nominal supere la tasa de interés,
porque el déficit primario es mayor que el alivio de $\gamma>i$.
\end{dato}

\begin{alerta}
El sitio presenta las cifras fiscales en metodología de caja. Las cifras oficiales en metodología de causación del
Ministerio de Hacienda pueden diferir. La regla fiscal de Colombia (Ley 1473 de 2011 y Ley 2155 de 2021) usa sus
propias definiciones.
\end{alerta}

% ---------------------------------------------------------------------
\chapter{Revisiones y datos provisionales}
\label{es:f:revisiones}

\section{Por qué cambian los datos publicados}
Las cifras oficiales se publican primero como \textbf{provisionales} y se revisan cuando llega información más
completa:
\begin{itemize}
\item el PIB trimestral y el ISE se revisan en cada publicación (nueva información de encuestas y registros, nuevos
  factores estacionales);
\item el PIB anual y departamental pasa por versiones provisional, preliminar y definitiva;
\item las cifras de comercio exterior del último año son provisionales;
\item los cambios de año base o de metodología reescriben toda la historia.
\end{itemize}
Cada versión publicada de una serie es una \textbf{vintage}. El sitio guarda solo la vintage vigente: la historia que
muestra es la de la última publicación, no lo que se sabía en cada fecha.

\begin{dato}[Ejemplo ilustrativo]
El crecimiento de un trimestre se publica como 3,0\% y, tres meses después, con más información, se revisa a 3,4\%. Con
la serie revisada, la brecha del producto de ese trimestre y la fase del ciclo pueden cambiar.
\end{dato}

\begin{alerta}
Las revisiones afectan sobre todo a los últimos datos, que son justamente los que más interesan. Una lectura que depende
de una décima en el último trimestre es frágil. La brecha ``en tiempo real'' elimina el uso de información futura del
filtro, pero no las revisiones de los datos del DANE.
\end{alerta}

\section{Datos que se sustituyen o se excluyen}
\begin{itemize}
\item La última inflación anual se toma de la cifra oficial del DANE (el índice con dos decimales puede diferir en una
  centésima).
\item Las observaciones aisladas que se alejan de forma anómala de sus dos vecinos (saltos de un día que se revierten)
  se conservan en los datos originales, pero se excluyen de los cálculos derivados.
\item Los meses incompletos de los registros (por ejemplo, matrículas de empresas del mes en curso) se descartan.
\end{itemize}
